\documentclass[12pt]{article}
\usepackage{cocina1}

\begin{document}
\begin{ficha}{Por teléfono}{Sobre el papel}{3}

\begin{encargo}
Nicoleta llama con dos fracciones. El cortador no llega: se cortan \textbf{sobre el papel}.
\end{encargo}

\bloque{1}{Mira cómo se hace}{Sobre el papel}\quad{\small ¿$\frac{3}{4}$ o $\frac{2}{3}$?}

\vspace{1.5mm}
\paso{1}\textbf{Multiplica la primera por el de abajo de la otra.}\quad \cortado{3}{4}{3}

\vspace{1mm}
\paso{2}\textbf{Multiplica la segunda por el de abajo de la primera.}\quad \cortado{2}{3}{4}

\vspace{1mm}
\paso{3}\textbf{Compara los de arriba.}\ \ 9 es más que 8:\quad \fA{3}{4}\ $>$\ \fB{2}{3}

\vspace{1mm}
\hspace*{6mm}{\small Más rápido: \textbf{en cruz}.}\quad\cruzllena{3}{4}{2}{3}

\vspace{1mm}
\begin{voz}[nicoleta]{NICOLETA, AL TELÉFONO}
Multiplica \textbf{en cruz}: el de arriba de una por el de abajo de la otra.
Gana la fracción que tiene el producto \textbf{más grande}.
\end{voz}

\bloque{2}{Ahora tú}{Sobre el papel}\quad{\small multiplica en cruz y pon el signo}

\vspace{1mm}
\textbf{a)}\ \cruz{2}{5}{1}{3}\hfill $\dfrac{2}{5}$\ \signo\ $\dfrac{1}{3}$

\vspace{1mm}
\textbf{b)}\ \cruz{3}{7}{2}{5}\hfill $\dfrac{3}{7}$\ \signo\ $\dfrac{2}{5}$

\alto{en cruz es el de \textbf{arriba} de una por el de \textbf{abajo} de la
otra. Arriba por arriba no vale.}

\reto{¿Son la misma pizza $\frac{2}{3}$ y $\frac{6}{9}$? Usa la cruz.}

\comprueba{Cada código abre esa llamada.}{%
\qr{1}{j=telefono\&a=3/4\&b=2/3}\hfill\qr{2a}{j=telefono\&a=2/5\&b=1/3}\hfill
\qr{2b}{j=telefono\&a=3/7\&b=2/5}\hfill\qr{reto}{j=telefono\&a=2/3\&b=6/9}}

\end{ficha}

% ─────────────────────────────── REVERSO ──────────────────────────────────
\begin{dorso}{Por teléfono}{Más llamadas}{3}

\begin{minipage}[t]{.6\linewidth}\vspace{0pt}
\bloque{1}{Con la cruz}{Sobre el papel}

\vspace{2mm}
{\renewcommand{\arraystretch}{2.6}
\begin{tabular}{@{}l@{\hspace{4mm}}l@{}}
\textbf{a)}\ $\dfrac{4}{7}\ \signo\ \dfrac{3}{5}$ & \hueco[1cm]\quad \hueco[1cm]\\
\textbf{b)}\ $\dfrac{5}{8}\ \signo\ \dfrac{3}{5}$ & \hueco[1cm]\quad \hueco[1cm]\\
\textbf{c)}\ $\dfrac{7}{9}\ \signo\ \dfrac{4}{5}$ & \hueco[1cm]\quad \hueco[1cm]\\
\textbf{d)}\ $\dfrac{5}{6}\ \signo\ \dfrac{4}{5}$ & \hueco[1cm]\quad \hueco[1cm]\\
\end{tabular}}
\end{minipage}\hfill
\begin{minipage}[t]{.37\linewidth}\vspace{6mm}
\tablamult
\end{minipage}

\bloque{2}{¿La misma pizza?}{Sobre el papel}\quad{\small si los dos productos son iguales, sí}

\vspace{2mm}
{\renewcommand{\arraystretch}{2.4}
\begin{tabular}{@{}l@{\hspace{12mm}}l@{\hspace{12mm}}l@{}}
\textbf{a)}\ $\dfrac{2}{4}$ y $\dfrac{3}{6}$:\ \hueco[1.4cm] &
\textbf{b)}\ $\dfrac{3}{5}$ y $\dfrac{6}{10}$:\ \hueco[1.4cm] &
\textbf{c)}\ $\dfrac{2}{3}$ y $\dfrac{3}{4}$:\ \hueco[1.4cm]\\
\end{tabular}}

\bloque{3}{Un problema de la pastelería}{Sobre el papel}

\vspace{1.5mm}
Nicoleta usa $\frac{2}{3}$ de un paquete de harina y el pinche $\frac{3}{5}$ de
otro igual. ¿Quién usa más harina?\hfill\hueco[2cm]

\vspace{4mm}
\begin{semaforo}
\sfila{Multiplico en cruz}
\sfila{Pongo bien el signo}
\sfila{Sé si dos fracciones son la misma pizza}
\end{semaforo}

\comprueba{Cada código abre esa llamada.}{%
\qr{1a}{j=telefono\&a=4/7\&b=3/5}\hfill\qr{1b}{j=telefono\&a=5/8\&b=3/5}\hfill
\qr{1c}{j=telefono\&a=7/9\&b=4/5}\hfill\qr{1d}{j=telefono\&a=5/6\&b=4/5}\hfill
\qr{3}{j=telefono\&a=2/3\&b=3/5}}
\end{dorso}
\end{document}
